\renewcommand{\docshorttitle}{Royal Express · Memoria técnica}
\renewcommand{\doctitle}{Royal Express}
\renewcommand{\docsubtitle}{Memoria técnica de Royal Express · sistema de gestión logística para LogistiXpress}
\renewcommand{\docversion}{1.0.0}
\renewcommand{\docdate}{Noviembre de 2026}
\hypersetup{
  pdftitle={Royal Express · Memoria técnica},
  pdfauthor={Royal Express · Equipo de desarrollo},
  pdfsubject={Memoria técnica de Royal Express · sistema de gestión logística para LogistiXpress},
  pdfkeywords={MVC, CRUD, PostgreSQL, FastAPI, RBAC, rutas, picking}
}
\begin{titlepage}
  \thispagestyle{empty}
  \centering
  \vspace*{1.8cm}
  {\color{RoyalOrange}\rule{\linewidth}{2.4pt}}\\[0.45cm]
  {\footnotesize\bfseries\color{RoyalGray}LOGISTIXPRESS · CENTRO DE DISTRIBUCIÓN}\\[0.35cm]
  {\Huge\bfseries\color{RoyalNavy} Royal Express}\\[0.25cm]
  {\Large\color{RoyalNavy} Memoria técnica}\\[0.1cm]
  {\large\color{RoyalGray} Arquitectura, modelo de datos e implementación}\\[0.45cm]
  {\color{RoyalOrange}\rule{\linewidth}{1pt}}\\[1.4cm]
  \begin{minipage}{0.86\linewidth}
    \centering
    {\large\color{RoyalGray} Documento de cierre técnico: decisiones de arquitectura, MVC, CRUD,
    RBAC, algoritmos de rutas, recepción, picking/packing, validación y despliegue OCI.}
  \end{minipage}\\[2.2cm]
  \begin{tabular}{@{}ll@{}}
    \textbf{Producto} & Royal Express · LogistiXpress \\
    \textbf{Documento} & Memoria técnica \\
    \textbf{Versión} & \docversion \\
    \textbf{Fecha} & \docdate \\
    \textbf{Plataforma} & FastAPI · PostgreSQL 16 · Docker/Podman \\
    \textbf{Estado} & Implementación operativa · versión 1.0.0 \\
  \end{tabular}\\[2.4cm]
  \vfill
  {\footnotesize\color{RoyalGray} Documento de Royal Express para LogistiXpress.\\
  Uso interno del centro de distribución.}
\end{titlepage}
\tableofcontents
\newpage
